\begin{textsample}{POS Dim 1 – 2019 – Score 21.00 – 2019/t001223.txt}  \label{ex:f1_pos_058}
It may seem like we ' re all standing on solid earth right now, but we ' re not. \textbf{The} \textbf{rocks} and \textbf{the} dirt underneath us are crisscrossed by tiny little fractures and empty spaces.
And these empty spaces are filled with astronomical quantities of microbes, such as these ones. \textbf{The} deepest that we found microbes so far into \textbf{the} earth is five kilometers down.
So like, if you pointed yourself at \textbf{the} ground and took off running into \textbf{the} ground, you could run an entire 5K race and microbes would line your whole path.
So you may not have ever thought about these microbes that are deep inside earth ' s crust, but you probably thought about \textbf{the} microbes living in our guts.
If you add up \textbf{the} gut microbiomes of all \textbf{the} people and all \textbf{the} animals on \textbf{the} \textbf{planet}, collectively, this weighs about 100, 000 tons.
This is a huge biome that we carry in our bellies every single day.
We should all be proud.
But it pales in comparison to \textbf{the} number of microbes that are covering \textbf{the} entire \textbf{surface} of \textbf{the} earth, like in our soils, our rivers and our oceans.
Collectively, these weigh about two billion tons.
But it turns out that \textbf{the} majority of microbes on earth aren ' t even in oceans or our guts or sewage treatment plants.
Most of them are actually inside \textbf{the} earth ' s crust.
So collectively, these weigh 40 billion tons.
This is one of \textbf{the} biggest biomes on \textbf{the} \textbf{planet}, and we didn ' t even know it existed until a few decades ago.
So \textbf{the} possibilities for what life is like down there, or what it might do for humans, are limitless.
This is a map showing a red dot for every place where we ' ve gotten pretty good deep subsurface samples with modern microbiological methods, and you may be impressed that we ' re getting a pretty good global coverage, but actually, if you remember that these are \textbf{the} only places that we have samples from, it looks a little worse.
If we were all in an alien spaceship, trying to reconstruct a map of \textbf{the} globe from only these samples, we ' d never be able to do it.
So people sometimes say to me,"Yeah, there ' s a lot of microbes in \textbf{the} subsurface, but... aren ' t they just kind of dormant?"This is a good point. \textbf{Relative} to a ficus plant or \textbf{the} measles or my kid ' s guinea pigs, these microbes probably aren ' t doing much of anything at all.
We know that they have to be slow, because there ' s so many of them.
If they all started dividing at \textbf{the} rate of E. coli, then they would double \textbf{the} entire weight of \textbf{the} earth, \textbf{rocks} included, over a single night.
In fact, many of them probably haven ' t even undergone a single cell division since \textbf{the} time of ancient Egypt.
Which is just crazy.
Like, how do you wrap your head around things that are so long- lived?
But I thought of an analogy that I really love, but it ' s weird and it ' s complicated.
So I hope that you can all go there with me.
Alright, let ' s try it.
It ' s like trying to figure out \textbf{the} life cycle of a tree... if you only lived for a day.
So like if human life span was only a day, and we lived in winter, then you would go your entire life without ever seeing a tree with a leaf on it.
And there would be so many human generations that would pass by within a single winter that you may not even have access to a history book that says anything other than \textbf{the} fact that trees are always lifeless sticks that don ' t do anything.
Of course, this is ridiculous.
We know that trees are just waiting for summer so they can reactivate.
But if \textbf{the} human life span were significantly shorter than that of trees, we might be completely oblivious to this totally mundane fact.
So when we say that these deep subsurface microbes are just dormant, are we like people who die after a day, trying to figure out how trees work?
What if these deep subsurface organisms are just waiting for their version of summer, but our lives are too short for us to see it?
If you take E. coli and \textbf{seal} it up in a test tube, with no food or nutrients, and leave it there for months to years, most of \textbf{the} cells die off, of course, because they ' re starving.
But a few of \textbf{the} cells survive.
If you take these old surviving cells and compete them, also under starvation conditions, against a new, fast-growing culture of E. coli, \textbf{the} grizzled old tough guys beat out \textbf{the} squeaky clean upstarts every single time.
So this is evidence there ' s actually an \textbf{evolutionary} payoff to being extraordinarily slow.
So it ' s possible that maybe we should not equate being slow with being unimportant.
Maybe these out-of-sight, out-of-mind microbes could actually be helpful to humanity.
OK, so as far as we know, there are two ways to do subsurface living. \textbf{The} first is to wait for food to trickle down from \textbf{the} \textbf{surface} world, like trying to eat \textbf{the} leftovers of a picnic that happened 1, 000 years ago.
Which is a crazy way to live, but shockingly seems to work out for a lot of microbes in earth. \textbf{The} other possibility is for a microbe to just say,"Nah, I don ' t need \textbf{the} \textbf{surface} world.
I ' m good down here."For microbes that go this route, they have to get everything that they need in order to survive from inside \textbf{the} earth.
Some things are actually easier for them to get.
They ' re more abundant inside \textbf{the} earth, like water or nutrients, like nitrogen and \textbf{iron} and phosphorus, or places to live.
These are things that we literally kill each other to get ahold of up at \textbf{the} \textbf{surface} world.
But in \textbf{the} subsurface, \textbf{the} problem is finding enough energy.
Up at \textbf{the} \textbf{surface}, plants can chemically knit together \textbf{carbon} dioxide molecules into yummy sugars as fast as \textbf{the} sun ' s \textbf{photons} hit their leaves.
But in \textbf{the} subsurface, of course, there ' s no sunlight, so this ecosystem has to solve \textbf{the} problem of who is going to make \textbf{the} food for everybody else. \textbf{The} subsurface needs something that ' s like a plant but it breathes \textbf{rocks}.
Luckily, such a thing exists, and it ' s called a chemolithoautotroph.
Which is a microbe that uses chemicals —"chemo,"from \textbf{rocks} —"litho,"to make food —"autotroph."And they can do this with a ton of different elements.
They can do this with sulphur, \textbf{iron}, manganese, nitrogen, \textbf{carbon}, some of them can use pure electrons, straight up.
Like, if you cut \textbf{the} end off of an \textbf{electrical} cord, they could breathe it like a snorkel.
These chemolithoautotrophs take \textbf{the} energy that they get from these processes and use it to make food, like plants do.
But we know that plants do more than just make food.
They also make a waste product, oxygen, which we are 100 percent dependent upon.
But \textbf{the} waste product that these chemolithoautotrophs make is often in \textbf{the} form of minerals, like rust or pyrite, like fool ' s gold, or carminites, like limestone.
So what we have are microbes that are really, really slow, like \textbf{rocks}, that get their energy from \textbf{rocks}, that make as their waste product other \textbf{rocks}.
So am I talking about \textbf{biology}, or am I talking about geology?
This stuff really blurs \textbf{the} lines.
So if I ' m going to do this thing, and I ' m going to be a biologist who studies microbes that kind of act like \textbf{rocks}, then I should probably start studying geology.
And what ' s \textbf{the} coolest part of geology?
Volcanoes.
This is looking inside \textbf{the} crater of Po{\EmojiFont �}{\EmojiFont �}s Volcano in Costa Rica.
Many volcanoes on earth arise because an oceanic tectonic \textbf{plate} crashes into a continental \textbf{plate}.
As this oceanic \textbf{plate} subducts or gets moved underneath this continental \textbf{plate}, things like water and \textbf{carbon} dioxide and other materials get squeezed out of it, like ringing a \textbf{wet} washcloth.
So in this way, subduction zones are like portals into \textbf{the} deep earth, where materials are exchanged between \textbf{the} \textbf{surface} and \textbf{the} subsurface world.
So I was recently invited by some of my colleagues in Costa Rica to come and work with them on some of \textbf{the} volcanoes.
And of course I said yes, because, I mean, Costa Rica is beautiful, but also because it sits on top of one of these subduction zones.
We wanted to ask \textbf{the} very specific question : Why is it that \textbf{the} \textbf{carbon} dioxide that comes out of this deeply buried oceanic tectonic \textbf{plate} is only coming out of \textbf{the} volcanoes?
Why don ' t we see it distributed throughout \textbf{the} entire subduction zone?
Do \textbf{the} microbes have something to do with that?
So this is a picture of me inside Po{\EmojiFont �}{\EmojiFont �}s Volcano, along with my colleague Donato Giovannelli.
That lake that we ' re standing next to is made of pure battery acid.
I know this because we were measuring \textbf{the} pH when this picture was taken.
And at some point while we were working inside \textbf{the} crater, I turned to my Costa Rican colleague Carlos Ram{\EmojiFont �}{\EmojiFont �}rez and I said,"Alright, if this thing starts erupting right now, what ' s our exit strategy?"And he said,"Oh, yeah, great question, it ' s totally easy.
Just turn around and \textbf{enjoy} \textbf{the} view.""Because it will be your last."And it may sound like he was being overly dramatic, but 54 days after I was standing next to that lake, this happened.
Audience : Oh!
Freaking \textbf{terrifying}, right?
This was \textbf{the} biggest eruption this volcano had had in 60-some-odd years, and not long after this video ends, \textbf{the} camera that was taking \textbf{the} video is obliterated and \textbf{the} entire lake that we had been sampling vaporizes completely.
But I also want to be clear that we were pretty sure this was not going to happen on \textbf{the} day that we were actually in \textbf{the} volcano, because Costa Rica monitors its volcanoes very carefully through \textbf{the} OVSICORI Institute, and we had scientists from that institute with us on that day.
But \textbf{the} fact that it erupted illustrates perfectly that if you want to look for where \textbf{carbon} dioxide gas is coming out of this oceanic \textbf{plate}, then you should look no further than \textbf{the} volcanoes themselves.
But if you go to Costa Rica, you may notice that in addition to these volcanoes there are tons of cozy little hot \textbf{springs} all over \textbf{the} place.
Some of \textbf{the} water in these hot \textbf{springs} is actually bubbling up from this deeply buried oceanic \textbf{plate}.
And our hypothesis was that there should be \textbf{carbon} dioxide bubbling up with it, but something deep \textbf{underground} was filtering it out.
So we spent two weeks driving all around Costa Rica, sampling every hot \textbf{spring} we could find — it was awful, let me tell you.
And then we spent \textbf{the} next two years measuring and analyzing data.
And if you ' re not a scientist, I ' ll just let you know that \textbf{the} big discoveries don ' t really happen when you ' re at a beautiful hot \textbf{spring} or on a public stage ; they happen when you ' re hunched over a messy \textbf{computer} or you ' re troubleshooting a difficult instrument, or you ' re Skyping your colleagues because you are completely confused about your data.
Scientific discoveries, kind of like deep subsurface microbes, can be very, very slow.
But in our case, this really paid off this one time.
We discovered that literally tons of \textbf{carbon} dioxide were coming out of this deeply buried oceanic \textbf{plate}.
And \textbf{the} thing that was keeping them \textbf{underground} and keeping it from being released out into \textbf{the} atmosphere was that deep \textbf{underground}, underneath all \textbf{the} adorable sloths and toucans of Costa Rica, were chemolithoautotrophs.
These microbes and \textbf{the} chemical processes that were happening around them were converting this \textbf{carbon} dioxide into carbonate mineral and locking it up \textbf{underground}.
Which makes you wonder : If these subsurface processes are so good at sucking up all \textbf{the} \textbf{carbon} dioxide coming from below them, could they also help us with a little \textbf{carbon} problem we ' ve got going on up at \textbf{the} \textbf{surface}?
Humans are releasing enough \textbf{carbon} dioxide into our atmosphere that we are decreasing \textbf{the} ability of our \textbf{planet} to support life as we know it.
And scientists and engineers and entrepreneurs are working on methods to pull \textbf{carbon} dioxide out of these point sources, so that they ' re not released into \textbf{the} atmosphere.
And they need to put it somewhere.
So for this reason, we need to keep studying places where this \textbf{carbon} might be stored, possibly in \textbf{the} subsurface, to know what ' s going to happen to it when it goes there.
Will these deep subsurface microbes be a problem because they ' re too slow to actually keep anything down there?
Or will they be helpful because they ' ll help convert this stuff to solid carbonate minerals?
If we can make such a big breakthrough just from one study that we did in Costa Rica, then imagine what else is waiting to be discovered down there.
This new field of geo-bio-chemistry, or deep subsurface \textbf{biology}, or whatever you want to call it, is going to have huge implications, not just for mitigating climate change, but possibly for understanding how life and earth have coevolved, or finding new products that are useful for industrial or medical applications.
Maybe even predicting \textbf{earthquakes} or finding life outside our \textbf{planet}.
It could even help us understand \textbf{the} origin of life itself.
Fortunately, I don ' t have to do this by myself.
I have amazing colleagues all over \textbf{the} world who are cracking into \textbf{the} mysteries of this deep subsurface world.
And it may seem like life buried deep within \textbf{the} earth ' s crust is so far away from our daily experiences that it ' s kind of irrelevant.
But \textbf{the} truth is that this weird, slow life may actually have \textbf{the} answers to some of \textbf{the} greatest mysteries of life on earth.
Thank you.

% matched lemmas: biology, carbon, computer, earthquake, electrical, enjoy, evolutionary, iron, photon, planet, plate, relative, rock, seal, spring, surface, terrifying, the, underground, wet
\end{textsample}
